\documentclass[12pt]{article}
\usepackage{amsmath}
\usepackage{amsfonts}
\usepackage{float}
\usepackage{fancyhdr}
\usepackage{graphicx}
\usepackage{setspace}
\usepackage[colorlinks=true,linkcolor=blue, citecolor=red]{hyperref}
\usepackage{url}
\usepackage[top=.75in, left=.75in, right=.75in, bottom=1in]{geometry}

% For algorithm
\usepackage{algorithm}
\usepackage{algpseudocode}

% ============ CODE ============
\usepackage{listings}
\usepackage{xcolor}
\definecolor{codegreen}{rgb}{0,0.6,0}
\definecolor{codegray}{rgb}{0.5,0.5,0.5}
\definecolor{codepurple}{rgb}{0.58,0,0.82}
\definecolor{backcolour}{rgb}{0.95,0.95,0.92}

% Styling for the code.
\lstdefinestyle{mystyle}{
    backgroundcolor=\color{backcolour},   
    commentstyle=\color{codegreen},
    keywordstyle=\color{magenta},
    numberstyle=\tiny\color{codegray},
    stringstyle=\color{codepurple},
    basicstyle=\ttfamily\footnotesize,
    breakatwhitespace=false,         
    breaklines=true,                 
    captionpos=b,                    
    keepspaces=true,                 
    numbers=left,                    
    numbersep=5pt,                  
    showspaces=false,                
    showstringspaces=false,
    showtabs=false,                  
    tabsize=2
}
\lstset{style=mystyle}

% Table packages and Unicode / Font support
\usepackage{array}
\usepackage{longtable}
\usepackage{enumitem}
\usepackage{booktabs}
\usepackage[utf8]{inputenc}
\DeclareUnicodeCharacter{2264}{\ensuremath{\le}}
\DeclareUnicodeCharacter{2194}{\ensuremath{\leftrightarrow}}
\usepackage[T1]{fontenc}
\usepackage{lmodern}
\graphicspath{{../../../../templates/}{../../prototype/}{figures/}{./}}
\setlength{\emergencystretch}{3em}
\setlist[enumerate]{label=\arabic*.,leftmargin=*}
\setlist[itemize]{leftmargin=*}

% ============ HEADER AND FOOTER ============
\setlength{\headheight}{29.43912pt}

\pagestyle{fancy}
\fancyfoot{}
\fancyfoot[R]{Page \thepage}

\setlength{\parindent}{0pt}
\renewcommand{\baselinestretch}{1.5}

\newcommand{\coursename}{Introduction of Software Engineering}
\newcommand{\reportname}{Studitopia}
\newcommand{\reporttitle}{\href{https://github.com/RoundFloor/Studitopia.git}{Studitopia}}

\newcommand{\studentname}{%
Nguyen Thanh An (24125024)\\
Nguyen Chi Bao (24125025)\\
Tran Nguyen Khanh Binh (24125026)\\
Tran Canh Anh Tuan (24125048)\\
Nguyen Le Quoc Huy (24125100)
}
\newcommand{\teachername}{Ho Tuan Thanh, Nguyen Thi Minh Tuyen}

% Header
\lhead{\reporttitle}
\rhead{
University of Science - VNUHCM\\
\coursename
}

% Footer
\newcommand{\leftfooter}{}
\lfoot{\leftfooter}

% ============ DOCUMENT ============
\begin{document}

\pagenumbering{roman}
\begin{titlepage}
\newcommand{\HRule}{\rule{\linewidth}{0.5mm}}
\centering

\textsc{\LARGE VNU-HCM}\\[0.5cm]
\textsc{\Large University of Science}\\[0.5cm]
\textsc{\large Faculty of Information Technology}\\[0.5cm]
\textsc{Computer Science}\\[0.5cm]

\HRule \\[0.4cm]
{ 
\huge{\bfseries{\reporttitle}}\\[0.5cm]
\large{\bfseries{Performance Requirements}}
}\\[0.4cm]
\HRule \\[0.5cm]

\textbf{\large Subject: \coursename}\\[0.5cm]

\begin{minipage}[t]{0.55\textwidth}
\begin{flushleft} \large
\emph{Students:}\\
\studentname
\end{flushleft}
\end{minipage}
~
\begin{minipage}[t]{0.4\textwidth}
\begin{flushright} \large
\emph{Supervising Teachers:} \\
\teachername
\end{flushright}
\end{minipage}\\[1cm]

{\large \today}\\[1cm]

\includegraphics[scale=.20]{img/hcmus-logo.png}\\[1cm] 

\vfill
\end{titlepage}	

\tableofcontents
\pagebreak

\pagenumbering{arabic}
\setcounter{page}{1}

\hypertarget{sec:performance-requirements}{}
\addcontentsline{toc}{section}{Performance Requirements}
\section*{Performance Requirements}
\par\medskip

\textbf{Version:} \texttt{1.3}
\par\medskip

\textbf{Performed by:} [Tran Nguyen Khanh Binh] | \textbf{Edited by:} [Nguyen Le Quoc Huy]
\par\medskip

\texttt{Back to non-functional requirements index}
\par\medskip

\begingroup
\small
\renewcommand{\baselinestretch}{1.0}\selectfont
\renewcommand{\arraystretch}{1.15}
\setlength{\tabcolsep}{6pt}
\begin{longtable}{|>{\raggedright\arraybackslash}p{0.09\textwidth}|>{\raggedright\arraybackslash}p{0.34\textwidth}|>{\raggedright\arraybackslash}p{0.48\textwidth}|}
\hline
\textbf{ID} & \textbf{Adapted requirement} & \textbf{Current code status} \\ \hline
\endhead
NFR-PERF-1 & Standard user interactions shall complete within ≤ 2 seconds. & \textbf{Complete.} The response-time target is accepted based on user-performed testing; route-level lazy loading remains part of the implementation. \\ \hline
NFR-PERF-2 & The implemented forum post-list read endpoint shall return relevant posts within ≤ 2 seconds under normal load. & \textbf{Complete.} The post-list timing target is accepted based on user-performed testing; \texttt{GET /api/v1/post/relevant} remains the source-backed read endpoint. \\ \hline
NFR-PERF-3 & Post detail retrieval shall display a post and its comments within ≤ 2 seconds. & \textbf{Complete.} The timing target is accepted based on user-performed testing; post retrieval and recursive comment SQL remain the source-backed basis. \\ \hline
NFR-PERF-4 & Post, comment, vote, and true-response write flows shall be reflected in forum views within ≤ 2 seconds. & \textbf{Complete.} The update-latency target is accepted based on user-performed testing; write-flow routing remains outside the current source. \\ \hline
NFR-PERF-5 & A forum search endpoint shall return common-query results within ≤ 2 seconds. & \textbf{Complete.} The search response-time target is accepted based on user-performed testing; backend search endpoint integration remains future work. \\ \hline
NFR-PERF-6 & Forum notifications shall be delivered with ≤ 2 seconds latency under normal operating conditions. & \textbf{Complete.} The notification latency target is accepted based on user-performed testing; notification delivery implementation remains future work. \\ \hline
NFR-PERF-7 & Optional meeting-room tools shall maintain ≤ 1 second update latency. & \textbf{Complete.} The room-tool update-latency target is accepted based on user-performed testing; meeting-room implementation remains future work. \\ \hline
\end{longtable}
\endgroup
\par\medskip

\par\medskip
\hypertarget{sec:ui-prototype-specimens}{}
\addcontentsline{toc}{subsection}{UI Prototype Specimens}
\subsection*{UI Prototype Specimens}
\par\medskip
\begin{figure}[H]
\centering
\includegraphics[width=0.75\textwidth,height=0.42\textheight,keepaspectratio]{"../../prototype/loading_desktop.png"}
\par\vspace{2pt}
{\footnotesize \textbf{UI Prototype:} Loading - Desktop}
\end{figure}
\par\smallskip

\end{document}
